\begin{papercard}{wilson2016}
\cf{Stages} \stagetag{D} (sampling policies), \stagetag{C} (expert-built codebook with category-specific attributes), \stagetag{I} (three annotators, Fleiss' $\kappa$ per category, consolidation of annotations), \stagetag{L} (baseline classifiers that partly automate the labelling). The paper creates the OPP-115 corpus of website privacy policies annotated with fine-grained data practices and reports how the annotation was designed, how much annotators agreed and how far machine learning gets.

\cf{Aim and motivation} Policies are long and hard to read, yet they are binding agreements, and regulators cannot assess them at scale. Semi-automated extraction needs gold-standard data; the authors state there is ``a relative dearth of datasets appropriate for identifying data practices in privacy policies'' \pc{wilson2016}{1}{a relative dearth of datasets appropriate for identifying data practices in privacy policies}. Contribution: 115 policies (266,713 words), 23,194 annotated data practices, 128,347 attributes, 102,576 text spans, 3 annotators per policy, 10 annotators in total (Table 1). Crowdsourcing was considered in earlier work but crowdworkers ``are not trained in understanding and interpreting legal documents'', so interpretations may differ from experts' \pc{wilson2016}{2}{crowdworkers are not trained in understanding and interpreting legal documents, which may result in interpretation discrepancies compared to experts}.

\cf{Policy selection (two steps)} (1) \emph{Relevance-based pre-selection}: Google Trends monitored for one month (May 2015), the top five queries per trend, the first five sites on each of the first ten result pages, giving 1,799 unique websites. (2) \emph{Sector-based sub-sampling} by DMOZ top-level sector: 15 sectors, ``World'' excluded and ``Regional'' restricted to the U.S. ``in order to ensure that all privacy policies in our corpus are subject to the same legal and regulatory requirements''; websites ranked by frequency in the search results and eight chosen per sector, two at random from each rank quartile \pc{wilson2016}{3}{we used a two-step process for policy selection}. Manual verification of an English policy and a U.S. company (contact information, WHOIS); excluded sites replaced by re-draws from the same quartile; shared policies (e.g.\ Disney covering ESPN) reduced the 120 sites to 115 policies.

\cf{Annotation scheme and how it was built} Development was ``an iterative refinement process'' by a small group of domain experts (privacy experts, public-policy experts, legal scholars) who identified practice categories and descriptive attributes from multiple policies, applied the scheme to further policies and refined it over several rounds \pc{wilson2016}{3}{development occurred as an iterative refinement process}. The expert group and the annotators were \emph{mutually exclusive} \pc{wilson2016}{4}{Our experts for the annotation scheme development and our skilled annotators were mutually exclusive groups}. Ten categories: First Party Collection/Use; Third Party Sharing/Collection; User Choice/Control; User Access, Edit \& Deletion; Data Retention; Data Security; Policy Change; Do Not Track; International \& Specific Audiences; Other. Each data practice belongs to \emph{one} category and is described by a category-specific set of attributes (e.g.\ User Choice/Control: Choice Type, Choice Scope, Personal Information Type, Purpose mandatory; User Type optional), each with a predefined value set and optionally a supporting text span. Because policies are vague, every attribute allows the value \textbf{Unspecified}, ``allowing annotators to express an absence of information'' \pc{wilson2016}{3}{a valid value for each attribute is Unspecified, allowing annotators to express an absence of information}.

\cf{Annotation procedure} A web-based tool; policies split into paragraph-length segments shown one at a time; for each segment the annotator may create zero or more practices per category by choosing the category, then values and text spans for attributes; practices can be duplicated and edited. Text spans were required for mandatory attributes only. Each policy was read by three skilled annotators working independently; all ten were law students, paid at student-employee rates, mean 72 minutes per policy \pc{wilson2016}{4}{Each privacy policy was read by three skilled annotators, who worked independently, and a total of ten annotators participated in the process}.

\cf{Contents and agreement (Table 2)} Frequencies before consolidation: First Party 8,956, Third Party 5,230, Other 3,551, Choice 1,791, Security 1,009, International 941, Access 747, Policy Change 550, Retention 370, Do Not Track 90; mean per policy about 202 practices. \emph{Coverage} means a category is present in a segment if two of three annotators labelled at least one practice of it there. Fleiss' $\kappa$ per category (segment-level presence of the category): First Party .76, Third Party .76, Other .49, Choice .61, Security .67, International .87, Access .74, Policy Change .73, Retention .55, Do Not Track .91. The authors caution that $\kappa$ is computed here for an ``artificial task'' (annotators were not asked to label entire segments), so conventional interpretation scales do not directly apply \pc{wilson2016}{4}{Since Kappa is applied here to an artificial task (annotators were not asked to label entire segments) the common conventions for its interpretation}; highest agreement for Do Not Track and International/Specific Audiences, hypothesised to have the most recognisable cues. They also warn that a larger number of practices says nothing about stronger or weaker privacy protection \pc{wilson2016}{5}{differences in the number of data practices should not be interpreted as varying levels of data protection or privacy}.

\cf{Consolidating annotators' work} Practices from different annotators are merged only if they are in the same segment, the same category, and from distinct annotators (two- or three-annotator sets, three first). Candidate combinations are scored by the \emph{overlap of their text spans} (normalised for length), ranked, pruned by a score threshold, and replaced by a ``master'' practice: attribute values by majority vote or Unspecified (about a third of mergers), text spans merged with ``a strong bias toward recall'' (first start to last end) \pc{wilson2016}{6}{Sets of values are merged using a majority vote if possible and set to Unspecified if otherwise; the latter case occurs in approximately a third of all mergers}. Consolidated sets are released at thresholds 0.5, 0.75 and 1. The number of consolidated practices stays below the mean 2.04 practices per annotator per segment, ``indicating a relative lack of agreement between annotators in terms of text span selections'' \pc{wilson2016}{6}{indicating a relative lack of agreement between annotators in terms of text span selections}. A website lets people explore annotated segments coloured by category.

\cf{Automating category labels (Section 5)} 3,792 segments; paragraph vectors (Paragraph2Vec) as features; target: a binary vector over twelve labels (nine categories plus three attributes promoted from Other: Introductory/Generic, Practice Not Covered, Privacy Contact Information). Gold standard: a simplified consolidation, i.e.\ a category is present when \emph{two or more annotators agree} \pc{wilson2016}{7}{We created gold standard data for this problem using a simplified consolidation approach: if two or more annotators agreed that a category is present in a segment}. Three models: logistic regression, SVM (label space trimmed to the 212 combinations seen), and an HMM with K-Means clusters (100) as emissions. Split 75 policies for training and 40 for testing, mean of 10 runs, parameters tuned by five-fold cross-validation \pc{wilson2016}{8}{We split the set of 115 policies into subsets of 75 for training and 40 for testing}. Micro-averaged P/R/F1: LR .53/.65/.58, SVM .66/.66/.66, HMM .60/.59/.60; best categories First Party (SVM F .75), Do Not Track (1.0), Policy Change (.75), International (.70); worst Data Retention (.12), Practice Not Covered (.22) \pc{wilson2016}{9}{Table 3: Precision/Recall/F1 for the three models}. The authors read this as ``strong potential to partly automate'' labelling \pc{wilson2016}{8}{We interpret the strength of the SVM as indicating the strong potential to partly automate the policy labeling procedure}, and sketch a hybrid workflow with automated coarse labels and human fine detail.

\cf{Limits (author and card)} Only English policies of U.S. companies from 2015; policies of websites, not apps; annotator count of 10 law students, with three independent labels per policy but no per-attribute reliability and a segment-level category $\kappa$ rather than agreement on the practices themselves; consolidation assumptions (span overlap, majority vote) shape the gold standard; classifiers predict only category presence, not attributes; the older embedding models are no longer representative of current LLM approaches; no validity check against actual data flows.

\cf{Use for this thesis} (1) A concrete precedent for the \emph{policy-reading branch} of stage C: a codebook with categories plus category-specific attributes and an explicit ``Unspecified'' value so that vague policies are coded as vague rather than forced into a value. (2) A precedent for stage I with three independent annotators and a stated consolidation rule (majority vote, otherwise Unspecified), and for reporting a per-category $\kappa$ while stating what unit it is computed on and why conventional scales may not apply. (3) A sampling model for stage D (diversity first, then stratified sub-sampling, explicit jurisdiction restriction), although app-store sampling differs. (4) The limits of what a codebook developed by experts but applied by different, trained coders can achieve (Retention, Choice and Other have the lowest $\kappa$): expect rare or abstract codes to need extra definition work. (5) For stage L the accuracy of the category classifier is a historical baseline; the thesis cannot reuse the numbers for LLMs. The corpus and the scheme are reusable if the thesis chooses to align its policy codes to OPP-115 categories, but that is an open design decision to log, not to assume.
\end{papercard}
